\documentclass[10pt,a4paper]{amsart}
\usepackage[margin=19mm]{geometry}
\usepackage{amsmath,amssymb,mathtools}
\usepackage[scr=rsfs,cal=euler]{mathalfa}
\usepackage{xfrac}
\usepackage[unicode,hidelinks,bookmarks=false]{hyperref}
\newcommand{\ie}{i.e.\ }
\newcommand{\cf}{cf.\ }
\newcommand{\resp}{respectively\ }
\newcommand{\Subsection}{\subsection}
\providecommand{\nobreakdash}{\nobreak\hbox{-}}
\setlength{\emergencystretch}{2em}
\multlinegap0pt
\usepackage{amssymb}
\usepackage{bm}
\usepackage[shortlabels]{enumitem}
\usepackage{algorithm}\usepackage{algpseudocode}
\usepackage{xcolor}
\usepackage{mathtools}
\newtheorem{theorem}{Theorem}
\title{\vspace{-10mm}\textbf{RAMPAGE: RAndomized Mid-Point for debiAsed Gradient Extrapolation}}
\author{Zhankun Luo, M. Berk Sahin, Antesh Upadhyay, Behzad Sharif, Abolfazl Hashemi}
\date{Source: arXiv:2603.22155v2 (CC BY 4.0)}
\begin{document}
\maketitle

\noindent\emph{Source and licence.} \vspace{-10mm}\textbf{RAMPAGE: RAndomized Mid-Point for debiAsed Gradient Extrapolation}. Version arXiv:2603.22155v2 (CC BY 4.0), distributed by arXiv under the Creative Commons Attribution 4.0 International licence, http://creativecommons.org/licenses/by/4.0/, as recorded in the arXiv licence metadata for this exact version and shown on the abstract page https://arxiv.org/abs/2603.22155v2. Full licence notice: \texttt{licenses/arxiv-2603.22155-CC-BY-4.0.txt}.\par\smallskip

\noindent\textit{Editorial scope.} Counted unit: the article's theorem thm:es\_ramp\_general (source0.tex lines 283--292). The article's complete original proof of it is delivered with it (source0.tex lines 807--879, one \texttt{proof} environment of 73 lines, which the article places away from the statement). No mathematics is written, repaired or re-derived: the statement and the proof are the article's own lines, in the article's own notation, and no retained line is substituted or re-flowed.  Retained uncounted as the source-local context the proof needs: lines 258--266.  Nothing was omitted from any retained span, so the package carries no omission record.  No retained line contains a citation, so the wrapper carries no bibliography and no bibliography entry.  No retained line contains a figure environment, an image-inclusion command or drawing code, so no image is delivered and the package declares no asset dependency.  Editorial formatting only, in addition: an AMS \texttt{amsart} preamble that restates the article's own declarations and macros used by the retained lines, namely the environments \texttt{theorem} on separate counters, together with the standard packages those lines need.  The retained environments print in this document as Theorem 1 in order of appearance, whereas the article prints the same items with the numbering of its own full document.  The complete native source of the article as distributed in the arXiv e-print for this exact version is bundled unchanged as \texttt{sources/197001/source0.tex}.
\begin{equation} \label{eq:rampage+}
\tag{RAMPAGE+}
\begin{aligned}
y_t &= \theta_t - 2\eta u_t F(\theta_t), \qquad
\tilde{y}_t = \theta_t - 2\eta \tilde{u}_t F(\theta_t), \\
\bar{F}_t &= \frac{1}{2} \big( F(y_t) + F(\tilde{y}_t) \big), \qquad
\theta_{t+1} = \theta_t - \eta \bar{F}_t.
\end{aligned}
\end{equation}

\begin{theorem} \label{thm:es_ramp_general}
Let $F: \mathbb{R}^p \to \mathbb{R}^p$ be an $L$-Lipschitz continuous and $\rho$-co-hypomonotone operator with $\rho \ge 0$. Assume $\mathrm{zer}(F) \neq \emptyset$. Let $\{\theta_t\}_{t \ge 0}$ denote the sequence generated by \eqref{eq:rampage+}. If the step size $\eta$ is chosen such that
\begin{equation}
C_{RA+}^{\rho,\eta} := \frac{3}{4} - \frac{2\rho}{\eta} - 16L^2 \left(\rho + \frac{\eta}{2}\right)^2 - 4L^2\eta \left(\rho + \frac{\eta}{2}\right) > 0,
\end{equation}
then for any $\theta^* \in \mathrm{zer}(F)$,
\begin{equation}
\min_{0 \le l \le k} \|F(\theta_l)\|^2 \le \frac{1}{k+1} \sum_{l=0}^k \|F(\theta_l)\|^2 \le \frac{\|\theta_0 - \theta^*\|^2}{C_{RA+}^{\rho,\eta} (k+1)}.
\end{equation}
\end{theorem}

\begin{proof}[Proof of Theorem \ref{thm:es_ramp_general}]
We track $\|\theta_{t+1} - \theta^*\|^2$. By the primary update rule $\theta_{t+1} = \theta_t - \eta \bar{F}_t$,
\begin{equation} \label{eq:esramp_base}
\|\theta_{t+1} - \theta^*\|^2 = \|\theta_t - \theta^*\|^2 - 2\eta \langle \bar{F}_t, \theta_t - \theta^* \rangle + \eta^2 \|\bar{F}_t\|^2.
\end{equation}
We decouple the inner product via $y_t$ and $\tilde{y}_t$
\begin{equation}
\begin{aligned}
-2\eta \langle \bar{F}_t, \theta_t - \theta^* \rangle &= -\eta \langle F(y_t), y_t - \theta^* \rangle - \eta \langle F(\tilde{y}_t), \tilde{y}_t - \theta^* \rangle \\
&\quad - \eta \langle F(y_t), \theta_t - y_t \rangle - \eta \langle F(\tilde{y}_t), \theta_t - \tilde{y}_t \rangle.
\end{aligned}
\end{equation}
Substituting $\theta_t - y_t = 2\eta u_t F(\theta_t)$ and $\theta_t - \tilde{y}_t = 2\eta \tilde{u}_t F(\theta_t)$, and applying the polarization identity $-2\langle a, b \rangle = \|a-b\|^2 - \|a\|^2 - \|b\|^2$ to the cross-terms yields
\begin{equation}
\begin{aligned}
-2\eta^2 u_t \langle F(y_t), F(\theta_t) \rangle &= -\eta^2 u_t \|F(y_t)\|^2 - \eta^2 u_t \|F(\theta_t)\|^2 + \eta^2 u_t \|F(y_t) - F(\theta_t)\|^2, \\
-2\eta^2 \tilde{u}_t \langle F(\tilde{y}_t), F(\theta_t) \rangle &= -\eta^2 \tilde{u}_t \|F(\tilde{y}_t)\|^2 - \eta^2 \tilde{u}_t \|F(\theta_t)\|^2 + \eta^2 \tilde{u}_t \|F(\tilde{y}_t) - F(\theta_t)\|^2.
\end{aligned}
\end{equation}
Since $u_t + \tilde{u}_t = 1$, summing these yields $-\eta^2 \|F(\theta_t)\|^2$. Applying the parallelogram law $\|a-b\|^2+\|a+b\|^2 = 2 \|a\|^2 + 2 \|b\|^2$ to $\eta^2 \|\bar{F}_t\|^2$ results in
\begin{equation}
\eta^2 \|\bar{F}_t\|^2 = \frac{\eta^2}{2}\|F(y_t)\|^2 + \frac{\eta^2}{2}\|F(\tilde{y}_t)\|^2 - \frac{\eta^2}{4}\|F(y_t) - F(\tilde{y}_t)\|^2.
\end{equation}
Define $A = F(y_t) - F(\theta_t)$ and $B = F(\tilde{y}_t) - F(\theta_t)$. Substituting these into \eqref{eq:esramp_base} results in
\begin{equation} \label{eq:esramp_grouped3}
\begin{aligned}
\|\theta_{t+1} - \theta^*\|^2 &= \|\theta_t - \theta^*\|^2 - \eta^2 \|F(\theta_t)\|^2 - \eta \langle F(y_t), y_t - \theta^* \rangle - \eta \langle F(\tilde{y}_t), \tilde{y}_t - \theta^* \rangle \\
&\quad + \eta^2 \left(\frac{1}{2} - u_t\right)\|F(y_t)\|^2 + \eta^2 \left(\frac{1}{2} - \tilde{u}_t\right)\|F(\tilde{y}_t)\|^2 \\
&\quad + \eta^2 u_t \|A\|^2 + \eta^2 \tilde{u}_t \|B\|^2 - \frac{\eta^2}{4} \|A - B\|^2.
\end{aligned}
\end{equation}
By $\rho$-co-hypomonotonicity of $F$, $-\eta \langle F(y), y - \theta^* \rangle \le \eta\rho \|F(y)\|^2$. Substituting this into \eqref{eq:esramp_grouped3} and defining $c_t = \frac{1}{2} - u_t$, such that $\frac{1}{2} - \tilde{u}_t = -c_t$, the coefficients for $\|F(y_t)\|^2$ and $\|F(\tilde{y}_t)\|^2$ become $\eta\rho + \eta^2 c_t$ and $\eta\rho - \eta^2 c_t$, respectively.

Expanding 
\begin{equation}
 \|F(y_t)\|^2 = \|F(\theta_t)\|^2 + 2\langle F(\theta_t), A \rangle + \|A\|^2
\end{equation}
and
\begin{equation}
 \|F(\tilde{y}_t)\|^2 = \|F(\theta_t)\|^2 + 2\langle F(\theta_t), B \rangle + \|B\|^2,
\end{equation}
the coefficient of $\|F(\theta_t)\|^2$ becomes 
\begin{equation}
 (\eta\rho + \eta^2 c_t) + (\eta\rho - \eta^2 c_t) = 2\eta\rho,
\end{equation}
The total coefficient for $\|F(\theta_t)\|^2$ is $-( \eta^2 - 2\eta\rho )$.
The cross-term is $2\eta \langle F(\theta_t), (\rho + \eta c_t)A + (\rho - \eta c_t)B \rangle$. Define $Y_t = \rho(A+B) + \eta c_t (A-B)$. Applying Young's inequality,
\begin{equation}
 2\eta \langle F(\theta_t), Y_t \rangle \le \frac{1}{4} \eta^2 \|F(\theta_t)\|^2 + 4 \|Y_t\|^2. 
\end{equation}
The remaining coefficient for $\|F(\theta_t)\|^2$ reduces to $-\frac{3}{4}\eta^2 + 2\eta\rho$. Since $|c_t| \le 1/2$, 
\begin{equation}
 \|Y_t\| \le (\rho + \frac{\eta}{2})(\|A\| + \|B\|).
\end{equation}

By $L$-Lipschitz continuity of $F$, $\|A\| \le 2L\eta u_t \|F(\theta_t)\|$ and $\|B\| \le 2L\eta \tilde{u}_t \|F(\theta_t)\|$. Since $u_t + \tilde{u}_t = 1$, $\|A\| + \|B\| \le 2L\eta \|F(\theta_t)\|$. Thus, $4\|Y_t\|^2 \le 16L^2\eta^2(\rho + \frac{\eta}{2})^2 \|F(\theta_t)\|^2$.

The quadratic terms sum to 
\begin{equation}
 V_Q = (\eta\rho + \eta^2 c_t)\|A\|^2 + (\eta\rho - \eta^2 c_t)\|B\|^2 + \eta^2 u_t \|A\|^2 + \eta^2 \tilde{u}_t \|B\|^2.
\end{equation}
Substituting $c_t = 1/2 - u_t$, the coefficients for both $\|A\|^2$ and $\|B\|^2$ evaluate to $\eta\rho + \frac{\eta^2}{2}$. Thus, $V_Q = \eta(\rho + \frac{\eta}{2})(\|A\|^2 + \|B\|^2)$. Using the Lipschitz bounds, 
\begin{equation}
 \|A\|^2 + \|B\|^2 \le 4L^2\eta^2 (u_t^2 + \tilde{u}_t^2) \|F(\theta_t)\|^2.
\end{equation}
Since $u_t^2 + \tilde{u}_t^2 \le 1$, we have $V_Q \le 4L^2\eta^3(\rho + \frac{\eta}{2}) \|F(\theta_t)\|^2$.

Substituting these bounds yields
\begin{equation}
\|\theta_{t+1} - \theta^*\|^2 \le \|\theta_t - \theta^*\|^2 - \eta^2 \left( \frac{3}{4} - \frac{2\rho}{\eta} - 16L^2\left(\rho + \frac{\eta}{2}\right)^2 - 4L^2\eta\left(\rho + \frac{\eta}{2}\right) \right) \|F(\theta_t)\|^2.
\end{equation}
The condition on $\eta$ ensures $C_{RA+}^{\rho,\eta} > 0$. Summing from $t=0$ to $k$, dividing by $(k+1)$, and evaluating the minimum completes the proof.
\end{proof}

\end{document}
